\section{Experiments}
\label{sec:experiments}

All experiments run in the MuJoCo digital twin (\SI{500}{Hz} physics,
\SI{50}{Hz} control) with scripted mock-hand trajectories that reproduce the
clutch semantics of Eq.~\eqref{eq:clutch}: the trajectory is defined by a
delta curve applied to the EE pose latched at ``grip press.'' Unless noted,
metrics aggregate both arms.

\paragraph{Metrics.}
Position error (\(\mathrm{mm}\), target vs.\ measured EE, mean/p95/max) and
its IK-only counterpart (target vs.\ commanded-configuration FK);
orientation error as the geodesic angle between target and measured EE
rotations (\(^\circ\)); \emph{roll lag} (\(\mathrm{ms}\)) as the
cross-correlation peak between commanded and measured roll about the tip
axis; command jerk RMS (\(\mathrm{rad/s^3}\)); peak joint velocity; minimum
joint-limit margin; and solver time. The out-of-envelope suite adds:
time outside, tracking error against the \emph{emitted} (post-policy)
target and against the \emph{raw} operator target while outside, peak joint
velocity while outside, and the re-entry recovery time (raw-target error
\(<\SI{10}{mm}\) sustained \SI{0.2}{s}).

\subsection{Tracking suite}
\label{sec:exp-tracking}

Three circle radii (3/5/\SI{8}{cm}) and four line-stroke directions on the
table plane, orientation held. Results in
Table~\ref{tab:tracking} (see also the per-trajectory tables in
\texttt{tables/}). The headline finding of the earlier benchmark holds:
full-cost 6D tracking (\texttt{pink\_full}) pays a large position penalty
on a 5-DoF arm (\SI{34.1}{mm} mean across the full suite vs.\
\SI{15.3}{mm} for \texttt{pink\_relaxed}), while \texttt{scipy\_ls}
(\SI{11.2}{mm}) and the new \texttt{telegrip} split IK (\SI{11.7}{mm})
lead the sweep. \texttt{telegrip}'s position subproblem is 3-DoF and
cannot be disturbed by orientation terms, giving it sub-millimeter
commanded-FK error at a solver cost 13\(\times\) below
\texttt{scipy\_ls} (\SI{0.19}{ms} vs.\ \SI{2.5}{ms} per tick).

\begin{table}[h]
  \centering
  \caption{Tracking suite, representative trajectory
  (\texttt{circle\_r5cm}). Full sweep in \texttt{tables/}.}
  \label{tab:tracking}
  \small
  \begin{tabular}{lrrrrrrr}
\toprule
method & $\bar{e}_{pos}$ (mm) & $e_{pos}^{95}$ (mm) & $\bar{e}_{ori}$ ($^\circ$) & lag (ms) & jerk (rad/s$^3$) & $\dot{q}_{max}$ & solve (ms) \\
\midrule
pink\_full & 33.7 & 51.7 & 1.99 & -- & 0.33 & 0.48 & 0.22 \\
pink\_relaxed & 17.4 & 23.1 & 9.22 & -- & 0.35 & 0.36 & 0.22 \\
dls & 24.8 & 38.6 & 2.45 & -- & 0.61 & 0.58 & 0.12 \\
mink & 28.3 & 44.3 & 2.01 & -- & 0.68 & 0.60 & 0.19 \\
scipy\_ls & 9.55 & 11.4 & 5.61 & -- & 0.76 & 0.55 & 2.54 \\
telegrip & 9.56 & 11.7 & 7.35 & -- & 0.71 & 0.57 & 0.18 \\
\bottomrule
\end{tabular}

\end{table}

\subsection{Wrist-agility suite}
\label{sec:exp-wrist}

Position held at the latch point while the target orientation oscillates:
roll (\(\pm45^\circ\)) and flex (\(\pm30^\circ\)) at 0.5/1/\SI{2}{Hz}, plus
a combined slow-circle + \SI{1}{Hz}-roll trajectory. This suite quantifies
the ``wrist is not agile'' complaint directly.
Table~\ref{tab:wrist} shows the \SI{1}{Hz} roll oscillation: the production
QP tracker lags the roll by \(\approx\)\SI{300}{ms} (about a third of the
oscillation period) with \(28^\circ\) mean orientation error, while the
Telegrip split IK cuts the lag to \(\approx\)\SI{120}{ms} and the error to
\(19^\circ\)~--- at the price of an order of magnitude more command jerk,
as the direct wrist mapping slams into the \SI{3}{rad/s} rate limit that
the \(\pm45^\circ\)@\SI{1}{Hz} profile (peak \SI{4.9}{rad/s}) exceeds by
design.

\begin{table}[h]
  \centering
  \caption{Wrist-agility suite, \SI{1}{Hz} roll oscillation
  (\texttt{wrist\_roll\_f1hz}).}
  \label{tab:wrist}
  \small
  \begin{tabular}{lrrrrrrr}
\toprule
method & $\bar{e}_{pos}$ (mm) & $e_{pos}^{95}$ (mm) & $\bar{e}_{ori}$ ($^\circ$) & lag (ms) & jerk (rad/s$^3$) & $\dot{q}_{max}$ & solve (ms) \\
\midrule
pink\_full & 10.7 & 14.7 & 16.9 & 100.0 & 58.1 & 4.72 & 0.21 \\
pink\_relaxed & 9.84 & 10.8 & 27.8 & 300.0 & 7.40 & 0.61 & 0.22 \\
dls & 9.91 & 11.8 & 22.0 & 140.0 & 211.1 & 3.00 & 0.11 \\
mink & 24.7 & 44.7 & 20.7 & 140.0 & 241.2 & 3.00 & 0.21 \\
scipy\_ls & 9.91 & 10.3 & 11.9 & 80.0 & 61.1 & 4.92 & 2.45 \\
telegrip & 9.91 & 10.7 & 19.5 & 120.0 & 373.3 & 3.00 & 0.19 \\
\bottomrule
\end{tabular}

\end{table}

\subsection{Out-of-envelope suite}
\label{sec:exp-envelope}

Three excursions: \texttt{envelope\_radial} (forward stroke \SI{0.30}{m},
\SI{2}{s} dwell far past \(r_{\max}\), smooth return),
\texttt{envelope\_swoop} (\SI{0.20}{m} dip below the floor), and
\texttt{envelope\_slide} (exit radially, translate laterally while outside,
return)~--- the last distinguishes policies that keep tracking the lateral
component (\texttt{project}, \texttt{slow}) from \texttt{freeze}, which
holds a fixed anchor. Representative results in
Table~\ref{tab:envelope}. With the legacy \texttt{warn} behavior the arm
grinds at its limits \(\approx\)\SI{131}{mm} from its own commanded target;
every active policy keeps the emitted-target error at
\(\approx\)\SI{26}{mm} with bounded joint velocity and \(\le\)\SI{1.8}{s}
recovery. \texttt{project} is the recommended default for data collection:
it minimizes raw-target deviation on the slide trajectory and re-enters
seamlessly.

\begin{table}[h]
  \centering
  \caption{Out-of-envelope suite, \texttt{envelope\_slide} with
  \texttt{pink\_relaxed} (policy comparison; full sweep in
  \texttt{tables/}).}
  \label{tab:envelope}
  \small
  \begin{tabular}{lrrrrrr}
\toprule
policy & $t_{out}$ (s) & $\bar{e}_{emit}$ (mm) & $\bar{e}_{raw}$ (mm) & $\dot{q}_{max}^{out}$ & $t_{rec}$ (s) & jerk (rad/s$^3$) \\
\midrule
freeze & 6.16 & 24.9 & 156.7 & 0.82 & -- & 3.03 \\
project & 6.16 & 26.9 & 147.7 & 0.87 & -- & 2.99 \\
slow & 6.16 & 26.8 & 148.0 & 0.77 & -- & 2.07 \\
warn & 6.16 & 129.9 & 129.9 & 0.99 & -- & 1.95 \\
\bottomrule
\end{tabular}

\end{table}

\subsection{End-to-end rehearsal}

The full production stack (One-Euro filter, clutch, armplane mapping,
envelope policy, IK thread) is exercised against the digital twin with the
scripted mock device: circles (\SI{11.6}{mm} mean EE tracking with
\texttt{pink\_relaxed}, \SI{7.0}{mm} with \texttt{telegrip}), wrist
oscillation, and envelope excursions with each policy. These runs validate
the wiring rather than the methods; the suites above are the controlled
comparison.
